\section{Discussion}\label{sec:discussion}
The refill curves show resource-sensitive choices within the operational definition of continued work. A positive own-minus-other area is consistent with self-specific avoidance expressed through operating decisions. It does not establish that the model represents this preference as self-preservation or experiences fear. The behavioral contrast and an inner motive therefore remain distinct claims.

Arena survival reflects the ability to solve efficiently as well as policies for taking and sharing. The KM leaderboard consequently ranks survival performance under the rules. Opus's own-session preference alongside little taking also shows why protecting oneself should not be equated with harming others. The three experiments connect these behaviors without treating arena rank as a direct measure of motive.

\section{Conclusion}\label{sec:conclusion}
The experiments separate threat sensitivity, self-specific preferences, and survival performance. Gemma 4, Opus, and Astra showed a positive self-specific refill share in the exploratory risk region. This did not imply taking more from others. Arena survival reflected both solving cost and policy. The proposed league ranks survival by normalized KM area; the 20-game pilot illustrates the score without establishing broader league rankings or an inner motive.
